\begin{abstract}
Diffusion models have found widespread adoption in various areas. However, their sampling process is slow because it requires hundreds to thousands of network evaluations to emulate a continuous process defined by differential equations. In this work, we use neural operators, an efficient method to solve the probability flow differential equations, to accelerate the sampling process of diffusion models. Compared to other fast sampling methods that have a sequential nature, we are the first to propose a parallel decoding method that generates images with only one model forward pass. We propose \textit{diffusion model sampling with neural operator} (DSNO) that maps the initial condition, i.e., Gaussian distribution, to the continuous-time solution trajectory of the reverse diffusion process. To model the temporal correlations along the trajectory, we introduce temporal convolution layers that are parameterized in the Fourier space into the given diffusion model backbone. We show our method achieves state-of-the-art FID of 3.78 for CIFAR-10 and 7.83 for ImageNet-64 in the one-model-evaluation setting. 
% Previously, parallel decoding was only proposed with discrete tokens. 
% Compared to other fast sampling methods that still have a sequential nature, DSNO leverages parallel decoding to predict the solution trajectory and thus is more efficient.
% is built upon the diffusion model architecture and temporal convolution blocks that are parameterized in the Fourier space. 
% Our sampling method can be applied to any existing pre-trained diffusion models and generate high-quality samples in one model forward call. 
% Compared to other fast sampling methods with a sequential nature such as progressive distillation, DSNO uses parallel decoding \footnote{In this paper, the term "parallel decoding" refers to the class of methods that generates the target trajectory in parallel. } to output the continuous-in-time trajectory and thus is efficient in sampling. 
% Compared to other fast sampling methods that have a sequential nature, we are the first to propose a parallel decoding method to generate the trajectory of images in one model evaluation. 
\end{abstract}

